% ============================================================
% Template base para artigos de survey/tutorial gerado pela
% skill survey-writing. Ajuste \documentclass conforme o venue
% confirmado na Fase 0 — classes comuns em Computação/Engenharia:
%   - IEEEtran           (periódicos/conferências IEEE)
%   - acmart             (periódicos/conferências ACM, ex.: ACM CSUR)
%   - sbc                (conferências/periódicos da SBC, ex.: SBRC, JBCS)
%   - article (fallback) usado abaixo — compila em qualquer
%     distribuição LaTeX padrão quando a classe do venue ainda
%     não estiver disponível no ambiente do usuário.
% Troque a linha abaixo pela classe exigida assim que o usuário
% confirmar o venue e disponibilizar o arquivo .cls, se necessário.
% ============================================================
\documentclass[11pt]{article}

\usepackage[utf8]{inputenc}
\usepackage[T1]{fontenc}
\usepackage{graphicx}
\usepackage{booktabs}
\usepackage{amsmath,amssymb}
\usepackage{hyperref}
\usepackage[numbers,sort&compress]{natbib}

\title{TÍTULO PROVISÓRIO — ver Fase 1 em research.md}
\author{AUTOR(ES)}
\date{\today}

\begin{document}
\maketitle

\begin{abstract}
% TODO (Fase 8): escrever por último, resumindo escopo, contribuições
% e principais achados/insights da síntese (não uma lista de tópicos).
\end{abstract}

% --------------------------------------------------------------
% Introdução: SEM subseções, SEM figuras/tabelas/gráficos (a
% primeira ilustração do artigo só aparece a partir da seção 2).
% Motivação, escopo (Fase 1), posicionamento SUPERFICIAL frente a
% surveys existentes (só o suficiente para situar a contribuição —
% a comparação detalhada e a eventual tabela comparativa vão para
% o Background ou para a subseção de taxonomia do corpo principal,
% nunca aqui), resumo em 1-2 frases do framework/taxonomia adotada
% (o detalhamento é exclusivo da primeira subseção do corpo
% principal), contribuições explícitas em bullets. Ver template
% próprio (assets/template-introducao.tex) e
% references/03-estrutura-taxonomias.md, seção "Regras de conteúdo
% da Introdução".
% --------------------------------------------------------------
\input{sections/01_introducao.tex}

% --------------------------------------------------------------
% Background/Preliminares (opcional — decidido na Fase 3). É aqui
% (ou, na ausência de Background, na subseção de taxonomia do
% corpo principal) que entram o aprofundamento do estado da arte,
% o detalhamento do framework/taxonomia e a eventual tabela
% comparativa com surveys relacionadas — nunca na Introdução. É
% também a partir desta seção (seção 2 do artigo) que figuras,
% tabelas e gráficos podem aparecer pela primeira vez.
% --------------------------------------------------------------
% \input{sections/02_background.tex}

% --------------------------------------------------------------
% CORPO PRINCIPAL — UMA ÚNICA seção (\section), cujo título é
% confirmado com o usuário na Fase 3 (ex.: "Panorama da Área",
% "Revisão da Literatura", ou um nome específico ao tema — nunca
% "Corpo Principal" literalmente no artigo final). Dentro dela, o
% arquivo abaixo contém, em sequência, uma \subsection para a
% taxonomia adotada — com o detalhamento completo do framework e,
% se não houver Background, também a tabela comparativa com
% surveys relacionadas —, uma \subsection por eixo principal da
% taxonomia, e uma \subsection por seção transversal pertinente
% (segurança/privacidade, datasets/benchmarks, padronização/
% indústria) — tudo no mesmo nível de subseção, no mesmo arquivo.
% Ver references/03-estrutura-taxonomias.md.
% --------------------------------------------------------------
\input{sections/03_corpo_principal.tex}

% --------------------------------------------------------------
% Seção de tutorial (apenas se Fase 5 aplicável) — seção própria,
% separada do corpo principal.
% --------------------------------------------------------------
% \input{sections/09_tutorial.tex}

% --------------------------------------------------------------
% Desafios e direções futuras (Fase 7) — seção dedicada, nunca
% espalhada pelo texto.
% --------------------------------------------------------------
\input{sections/10_desafios_futuros.tex}

% --------------------------------------------------------------
% Conclusão.
% --------------------------------------------------------------
\input{sections/11_conclusao.tex}

% Estilo de bibliografia depende do venue confirmado na Fase 0:
% IEEEtran (IEEE), ACM-Reference-Format (ACM), sbc (SBC), ou
% plainnat/abbrvnat como fallback genérico.
\bibliographystyle{IEEEtran}
\bibliography{refs}

\end{document}
